\input{template.tex}

\ifPDFTeX
  \renewcommand{\rmdefault}{cmr}
\fi

\def\SheetNo{8}

\begin{document}

\sheettitle{Графы. Деревья и обходы}

\begin{problems}

% Второй лист, задача 4.
\item В связном графе нет циклов длины менее $5$ и степени всех вершин не менее $10$. Докажите, что в нём не менее $101$ вершины.

% Первый лист, задача 5.
\item Кодовый замок имеет $7$ кнопок, код состоит из $4$ нажатий (кнопки могут повторяться). Замок помнит только последние $4$ нажатия и открывается, как только они совпадут с кодом. За какое наименьшее число нажатий можно гарантированно открыть замок, не зная кода?

% Второй лист, задача 2.
\item В марсианском метро $2025$ станций. От любой станции до любой другой можно проехать. Забастовочный комитет хочет закрыть проезд через одну из станций так, чтобы между всеми остальными станциями был возможен проезд. Докажите, что такая станция найдётся.

% Первый лист, задача 8.
\item Сколько различных гамильтоновых циклов в полном двудольном графе $K_{n,n}$? Циклы, отличающиеся только выбором начальной вершины или направлением обхода, считаются одинаковыми.

% Второй лист, задача 9.
\item Дано дерево. За ход разрешается стереть одно ребро, ведущее в висячую вершину, и соединить эту вершину ребром с любой другой вершиной. За какое наименьшее число ходов можно из одного заданного дерева на $n$ вершинах гарантированно получить другое заданное дерево на $n$ вершинах (все вершины считаются одинаковыми)?

% Первый лист, задача 1.
\item Пусть $T$ --- дерево с $n\ge2$ вершинами, а $\Delta$ --- максимальная степень вершины в $T$. Докажите, что в $T$ не меньше $\Delta$ листьев.

% Второй лист, задача 8.
\item Дан связный граф, у которого $200$ вершин нечётной степени, а остальные имеют чётную степень. Докажите, что вершины этого графа можно покрыть $100$ путями (не обязательно простыми!), не имеющими общих рёбер.

% Первый лист, задача 3.
\item Может ли шахматный конь обойти доску $5\times5$, побывав на каждой клетке ровно один раз, и последним ходом вернуться на исходную клетку? Сформулируйте ответ на языке гамильтоновых циклов.

% Второй лист, задача 6.
\item Докажите, что у дерева может быть не более двух центров.

% Первый лист, задача 6.
\item В стране $15$ городов, некоторые пары городов соединены авиалиниями, каждая из которых принадлежит одной из трёх авиакомпаний. Известно, что если любая одна компания прекратит полёты, всё равно из каждого города можно будет добраться в любой другой рейсами двух оставшихся компаний (возможно, с пересадками). Какое наименьшее число авиалиний может быть в стране?

% Второй лист, задача 3.
\item В связном графе нашлись три простых пути максимальной длины, каждые два из которых пересекаются ровно по одной вершине. Докажите, что у всех трёх путей есть общая вершина.

% Первый лист, задача 2.
\item
\begin{parts}
  \item Можно ли нарисовать все рёбра полного графа $K_5$, не отрывая карандаша от бумаги и не проводя ни одно ребро дважды, и вернуться в исходную точку?
  \item Тот же вопрос для $K_6$ (без требования вернуться).
  \item За какое наименьшее число «росчерков» (маршрутов без повторения рёбер, не имеющих общих рёбер между собой) можно нарисовать $K_6$?
\end{parts}

% Второй лист, задача 5. В пункте а исправлена формулировка для n=4:
% два отсекаемых треугольника могут соответствовать одной диагонали.
\item Выпуклый $n$-угольник разрезан непересекающимися диагоналями на треугольники.
\begin{parts}
  \item Докажите, что при $n\ge4$ найдутся два треугольника триангуляции, каждый из которых отсекается от многоугольника одной из проведённых диагоналей.
  \item Докажите, что при $n\ge6$ найдётся диагональ, отрезающая четырёхугольник или пятиугольник.
\end{parts}

% Первый лист, задача 7.
\item Можно ли расставить $777$ шахматных коней на доске $2017\times2017$ так, чтобы каждый из них бил ровно $4$ других коня?

% Второй лист, задача 1.
\item В стране любые два города связаны авиалинией одной из двух авиакомпаний. Докажите, что можно закрыть одну из авиакомпаний так, что по-прежнему от любого города можно будет добраться до любого другого.

% Первый лист, задача 4.
\item Связный простой граф $G$ таков, что степень каждой вершины равна $4$. Докажите, что рёбра $G$ можно покрасить в красный и синий цвета так, чтобы у каждой вершины было ровно два красных и два синих ребра.

% Второй лист, задача 7.
\item Пусть $G$ --- связный граф, причём отличный от дерева. Докажите, что $G$ имеет хотя бы три вершины, удаление любой из которых не нарушает связности.

\item В стране ровно $1000$ городов, некоторые пары городов соединены двусторонними авиалиниями. Известно, что для любого натурального $k\le500$ выполнено следующее утверждение: если выбрать любое множество $A$ из $k$ городов, то найдётся хотя бы $k$ городов, не принадлежащих $A$, каждый из которых соединён авиалинией хотя бы с одним городом из $A$. Какое наименьшее количество авиалиний может быть в этой стране?

\end{problems}

\end{document}
